\documentclass[12pt,a4paper]{article}
\usepackage[utf8]{inputenc}
\usepackage{amsmath, amssymb}
\usepackage{graphicx}
\usepackage{geometry}
\usepackage{setspace}
\usepackage{graphicx}
\usepackage{float}
\usepackage{subcaption}
\usepackage{epigraph}
\geometry{margin=1in}
\onehalfspacing
\title{Comparative Analysis of Statistical Time-Series Models and Machine Learning Approaches for Stock Price Prediction}
\author{Norbert Horváth}
\date{2026}
\begin{document}
\maketitle
\epigraph{
``I take the market-efficiency hypothesis to be the simple statement that security prices fully reflect all available information''
}{-- Eugene Fama}
\begin{abstract}
This paper explores the feasibility of predicting stock price movements using machine learning methods. It compares traditional statistical models with modern machine learning approaches and evaluates their ability to capture trends and generate actionable signals. The study emphasizes practical limitations such as market volatility, transaction costs, and nonlinear behavior of financial data.
\end{abstract}

\section{Introduction}
Financial markets are often characterized through the lens of the efficient market hypothesis (EMH), which asserts that security prices fully reflect all publicly available information \cite{7104d0bd-69aa-330c-badb-94aa0c8f5bac}.
Under a strict interpretation of the EMH, no model -- statistical or otherwise -- should be able to generate systematically profitable forecasts from past prices alone. Nevertheless, a substantial body of empirical work documents departures from strict efficiency, including volatility clustering, momentum effects, and short-horizon predictability, which motivates the continued study of forecasting models for financial time series \cite{Tsay_2010}.
This paper undertakes a comparative study of two broad forecasting paradigms applied to stock price data: classical statistical time-series models, which are linear, interpretable, and grounded in well-established asymptotic theory; and machine learning approaches, which are nonlinear, data-driven, and impose fewer distributional assumptions. Specifically, we examine autoregressive and moving-average models (AR, MA, ARIMA, SARIMA), long-memory models (ARFIMA), and conditional heteroskedasticity models (ARCH/GARCH), alongside regularized regression (ridge, lasso), support vector regression, random forests, and -- in later sections -- neural network architectures. The central question is not just which family of models achieves the lowest point-forecast error, but whether either paradigm can reliably identify the \emph{direction} of near-term price movements once realistic frictions are taken into account.
Even without precise price-level forecasts, correctly identifying the direction of a future move can, in principle, provide the basis for a profitable trading strategy. In practice, however, transaction costs, execution slippage, and short-term volatility can quickly erode small statistical edges, so that a model's forecast accuracy does not automatically translate into economic profitability \cite{49018680cf534023b6ae0f4e933c475d}.
This study therefore investigates whether machine learning models can outperform classical statistical benchmarks once these practical constraints are accounted for, and whether any performance gains are large enough to survive realistic trading frictions.

% =====================================================================
\section{Statistical Models}

Before introducing any model, it is worth confronting the problem that makes
financial forecasting unusual. Figure~\ref{fig:illusion} shows four price
series. Each one has stretches that a market commentator could describe
without hesitation: a rally that stalls and reverses, a long decline that
finds support, a sideways range that finally breaks upward. None of these
descriptions mean anything, because all four series were produced by
repeatedly tossing a fair coin and moving the price one unit up on heads and
one unit down on tails. There is no trend, no momentum and no support level
in any of them, because there is nothing in them at all except accumulated
noise.

\begin{figure}[ht]
\centering
\includegraphics[width=\linewidth]{figures/fig_rw_illusion.pdf}
\caption{Four simulated price paths, each generated by 500 tosses of a fair
coin. The apparent trends, reversals and consolidations are artefacts of
accumulated randomness. Any forecasting method must first demonstrate that it
is detecting something a coin toss does not already produce.}
\label{fig:illusion}
\end{figure}

This is the central difficulty. Human pattern recognition, and equally a
sufficiently flexible statistical or machine learning model, will find
structure in Figure~\ref{fig:illusion} whether or not any structure exists.
The purpose of the statistical models in this section is therefore twofold.
They provide forecasts, but more importantly they provide \emph{disciplined
tests} of whether an apparent pattern is real: each model comes with an
explicit description of what it assumes, a procedure for estimating its
parameters, and diagnostic tests that reveal when those assumptions fail.
This mathematical interpretability and theoretical foundation is what has
kept these models in use despite the volatility, noise and nonlinearity of
real markets \cite{49018680cf534023b6ae0f4e933c475d}.

Throughout this section the observed series is written $X_t$, where $t$
indexes trading days, and $r_t$ denotes a return computed from that series.
Every model is presented in the same way: the equation, a term-by-term
reading of what each symbol contributes, the conditions under which the model
is valid, how its parameters are estimated, and a figure showing what the
model actually detects when applied to data whose true structure is known.

\medskip
\noindent\textbf{A note on the figures.} All figures in this section are
generated from \emph{simulated} series with parameters chosen by the author.
This is deliberate. When data are simulated, the correct answer is known in
advance, so a figure can demonstrate that a diagnostic recovers a truth we
control rather than merely producing a plausible-looking picture. The
corresponding analysis on real market data is the subject of the empirical
section.

% ---------------------------------------------------------------------
\subsection{Implementation Framework and Methodological Constraints}

All models in this paper are implemented in Python. Classical time-series
estimation uses \texttt{statsmodels} \cite{Seabold_2010}, which provides the
ARIMA and SARIMA families, unit-root tests, the exponential smoothing state
space models and Markov-switching regressions. Automated order selection uses
\texttt{pmdarima}, conditional-volatility models use the \texttt{arch}
package, the regression and tree-based methods of the following section use
\texttt{scikit-learn}, and the neural network architectures introduced later
use \texttt{PyTorch}. Data handling uses \texttt{pandas} and \texttt{NumPy},
and all figures are produced with \texttt{matplotlib}. No component of the
study requires a compiled language: the series involved are short by modern
standards, typically a few thousand daily observations, and estimation time is
negligible compared with the effort of specifying and validating the models
correctly.

To keep comparisons between very different model families meaningful, the
following constraints apply uniformly to every model in this paper.

\begin{enumerate}

\item \textbf{Stationarity is established before estimation.} Models in the
ARIMA family assume the series has a stable mean and variance. Each series is
therefore tested with the Augmented Dickey--Fuller test \cite{Dickey_1979},
whose null hypothesis is that a unit root is present, and with the KPSS test,
whose null hypothesis is the opposite. Using both matters: failing to reject
the ADF null is not evidence of nonstationarity, only an absence of evidence
against it, so the two tests are read together and disagreement is reported
rather than resolved by preference.

\item \textbf{Model order is selected by information criteria, not by fit.}
Adding parameters always improves in-sample fit, so orders are chosen by
minimising the Akaike Information Criterion, which penalises complexity, and
cross-checked against the correlogram diagnostics described below.

\item \textbf{Residuals must be indistinguishable from noise.} A fitted model
is accepted only if what it failed to explain contains no remaining
structure. Residuals are tested for autocorrelation with the Ljung--Box test
\cite{Ljung_1978}, for remaining volatility structure with Engle's ARCH-LM
test, and for normality with the Jarque--Bera test.

\item \textbf{No information from the future may enter the past.} All splits
respect chronological order, and models are compared by rolling-origin
(walk-forward) re-estimation rather than random $k$-fold cross-validation. The
latter would train a model on observations that occur after the ones it is
asked to predict, which inflates measured accuracy in a way that cannot be
reproduced in live trading.

\item \textbf{All models are scored on the same terms.} Point-forecast
accuracy is measured by RMSE, MAE and MAPE; directional accuracy is measured
as the proportion of correctly predicted signs of change. Where a forecast is
converted into a trading signal, transaction costs and slippage are subtracted
before any claim of profitability is made.

\item \textbf{Every result must be reproducible.} Random seeds, package
versions and the exact data window are recorded for each experiment.

\end{enumerate}

These constraints are stated once here and assumed thereafter.

% ---------------------------------------------------------------------
\subsection{The Random Walk: the Benchmark That Must Be Beaten}

The random walk is the formal statement of the process that generated
Figure~\ref{fig:illusion}, and it is the direct testable implication of the
efficient market hypothesis discussed in Section~1. If prices already reflect
all available information, then the only reason the price changes is the
arrival of \emph{new} information, which by definition could not have been
known yesterday. The best available forecast of tomorrow's price is therefore
today's price, plus at most a constant long-run drift. The model is

\begin{equation}
X_t = \mu + X_{t-1} + \epsilon_t .
\end{equation}

Reading this term by term:
\begin{itemize}
\item $X_t$ is the price today and $X_{t-1}$ the price yesterday. The
      appearance of $X_{t-1}$ with a coefficient of exactly one is the whole
      content of the model: the entire past is compressed into the most recent
      observation, and nothing before yesterday carries any information.
\item $\mu$ is the drift, a constant added every period. With $\mu = 0$ the
      series has no systematic tendency to rise or fall; with $\mu > 0$ it
      tends to rise, which is the usual case for equity indices over long
      horizons.
\item $\epsilon_t$ is the shock arriving at time $t$, assumed to be white
      noise: mean zero, constant variance $\sigma^2$, and uncorrelated with
      shocks at every other date. It is the mathematical representation of
      genuinely new information.
\end{itemize}

Rearranging isolates what the model claims is predictable:
\begin{equation}
\Delta X_t \;=\; X_t - X_{t-1} \;=\; \mu + \epsilon_t .
\end{equation}
The change in price is a constant plus unforecastable noise. This is why the
random walk is such a demanding benchmark: it asserts that the only
predictable component of a price series is its average drift.

\paragraph{Why uncertainty grows with the horizon.} Iterating the model $h$
steps forward gives $X_{t+h} = X_t + h\mu + \sum_{i=1}^{h}\epsilon_{t+i}$.
Because the shocks are uncorrelated, their variances add:
\begin{equation}
\mathrm{Var}(X_{t+h} \mid X_t) \;=\; h\sigma^2 ,
\qquad
\mathrm{sd}(X_{t+h} \mid X_t) \;=\; \sigma\sqrt{h}.
\end{equation}
Uncertainty therefore grows with the square root of the horizon. The optimal
forecast and its interval are
\begin{equation}
\hat{X}_{t+h \mid t} = X_t + h\hat\mu ,
\qquad
\hat{X}_{t+h \mid t} \pm 1.96\,\hat\sigma\sqrt{h},
\end{equation}
shown in Figure~\ref{fig:rwfore}. The point forecast is a straight line and
the interval opens as a parabola. Note that the shocks never decay: a shock
today shifts the level of the series permanently, which is exactly what
distinguishes a random walk from the mean-reverting models that follow.

\begin{figure}[ht]
\centering
\includegraphics[width=\linewidth]{figures/fig_rw_forecast.pdf}
\caption{The random-walk forecast. Because every shock permanently shifts the
level and shock variances accumulate, the point forecast is a straight line
extrapolating the drift, and the interval widens in proportion to
$\sqrt{h}$.}
\label{fig:rwfore}
\end{figure}

\paragraph{Role in this study.} The random walk requires estimating a single
parameter, the mean change $\hat\mu$, and yet it is notoriously difficult to
beat out of sample on price levels at short horizons. It is therefore used
here not as a competitor of interest but as the no-skill benchmark. Any model,
statistical or machine learning, that does not improve on it under the
protocol of Section~2.1 provides no evidence against weak-form efficiency,
whatever its absolute error happens to be.

% ---------------------------------------------------------------------
\subsection{Autoregressive Model (AR)}

The random walk allows only one specific dependence on the past, a
coefficient of one on yesterday's value. The autoregressive model relaxes this
by letting the series depend on several past values with coefficients that are
estimated rather than assumed. In plain terms it describes a series with
memory that fades: a shock pushes the series away from its mean, and the
series then drifts back, with the speed of return governed by the
coefficients. The AR($p$) model is

\begin{equation}
X_t = \sum_{i=1}^{p} \phi_i X_{t-i} + \epsilon_t ,
\qquad \epsilon_t \sim \mathrm{WN}(0,\sigma^2).
\end{equation}

Term by term:
\begin{itemize}
\item $p$ is the order: how many past observations enter the model. An AR(1)
      looks back one day, an AR(5) looks back a full trading week.
\item $\phi_i$ is the weight on the observation $i$ periods ago. Its sign and
      size say how much of that day's deviation from the mean persists into
      today. A large positive $\phi_1$ means momentum at the daily scale; a
      negative $\phi_1$ means the series tends to reverse.
\item $\epsilon_t$ is again white noise, the part of today's movement that no
      combination of past values can account for.
\end{itemize}

\paragraph{Compact form and the stationarity condition.} Introducing the
backshift operator $B$, defined by $BX_t = X_{t-1}$, the model becomes
\begin{equation}
\phi(B)X_t = \epsilon_t ,
\qquad
\phi(B) = 1 - \phi_1 B - \phi_2 B^2 - \cdots - \phi_p B^p ,
\end{equation}
where $\phi(B)$ is called the autoregressive polynomial. The process is
stationary, meaning its mean and variance do not change over time, if and only
if every root $z$ of the equation
\begin{equation}
\phi(z) = 0
\end{equation}
satisfies $|z| > 1$, that is, lies outside the unit circle in the complex
plane. For the AR(1) case $X_t = \phi_1 X_{t-1} + \epsilon_t$ the single root
is $z = 1/\phi_1$, so the condition reduces to the familiar
\begin{equation}
|\phi_1| < 1 .
\end{equation}
The interpretation is direct. If $|\phi_1| < 1$, a shock is multiplied by
$\phi_1$ each period and dies out geometrically. If $\phi_1 = 1$ the shock
never decays and we are back to the random walk of Section~2.2, which is
precisely the boundary case this condition excludes. The random walk is thus
the AR(1) model with a unit root, and the ADF test of Section~2.1 is in
essence a test of whether $\phi_1$ equals one.

\paragraph{How the order is identified.} For a stationary AR(1) the
autocorrelation between observations $k$ days apart is
\begin{equation}
\rho_k = \phi_1^{\,k},
\end{equation}
which decays geometrically but never reaches exactly zero: because today
depends on yesterday and yesterday on the day before, today and the day before
are correlated indirectly. The partial autocorrelation removes those indirect
paths, measuring the correlation at lag $k$ after controlling for all shorter
lags. For an AR($p$) it is therefore zero beyond lag $p$. This gives the
diagnostic signature seen in the lower panels of Figure~\ref{fig:ar1}: a
gradually decaying ACF together with a PACF that cuts off sharply identifies
an AR process, and the lag at which the PACF cuts off estimates $p$
\cite{Box_2015}.

\begin{figure}[ht]
\centering
\includegraphics[width=\linewidth]{figures/fig_ar1.pdf}
\caption{A simulated AR(1) series with $\phi_1 = 0.7$ ($n = 400$). The ACF
decays geometrically as $0.7^{k}$, while the PACF is significant at lag 1 and
falls inside the 95\% white-noise band (shaded) thereafter, correctly
identifying the order $p = 1$.}
\label{fig:ar1}
\end{figure}

\paragraph{Estimation.} Multiplying the model by $X_{t-k}$ and taking
expectations yields the Yule--Walker equations, which relate the coefficients
to the autocorrelations of the series and can be solved directly:
\begin{equation}
\hat{\boldsymbol{\phi}} = \hat{R}_p^{-1} \hat{\boldsymbol{\rho}}_p ,
\end{equation}
where $\hat{\boldsymbol{\rho}}_p = (\hat\rho_1,\dots,\hat\rho_p)^{\!\top}$
collects the sample autocorrelations and $\hat{R}_p$ is the matrix whose
$(i,j)$ entry is $\hat\rho_{|i-j|}$. This estimator is consistent and requires
no numerical optimisation, but it is less efficient than maximum likelihood,
which is used as the default throughout this study.

\paragraph{Limitations.} An AR model is linear and has a single fixed
variance. It cannot represent the clustering of volatility addressed in
Section~2.10, nor abrupt changes in behaviour addressed in Section~2.11, nor
any nonlinear relationship of the kind the machine learning models of
Section~3 are designed to capture.

% ---------------------------------------------------------------------
\subsection{Moving Average Model (MA)}

The moving average model expresses the series not in terms of its own past
values but in terms of its own past \emph{errors}. In plain terms, it
describes a series that reacts to surprises for a fixed number of periods and
then forgets them completely -- a news item that moves the price today and
tomorrow, and has no effect at all on the third day. The MA($q$) model is

\begin{equation}
X_t = \epsilon_t + \sum_{i=1}^{q}\theta_i \epsilon_{t-i}.
\end{equation}

Term by term:
\begin{itemize}
\item $q$ is the order: how many past shocks still influence the present.
\item $\theta_i$ is the weight carried by the shock from $i$ periods ago. Once
      $i$ exceeds $q$ the weight is exactly zero, which is the defining
      feature of the model: memory ends abruptly rather than fading.
\item $\epsilon_t$ is today's shock, which enters with coefficient one.
\end{itemize}

Because the right-hand side is a finite sum of white-noise terms, an MA
process is stationary for \emph{any} values of $\theta_i$; no analogue of the
AR stationarity condition is needed.

\paragraph{Invertibility.} A separate condition is required instead. Writing
the model with the moving-average polynomial,
\begin{equation}
X_t = \theta(B)\epsilon_t ,
\qquad
\theta(B) = 1 + \theta_1 B + \cdots + \theta_q B^q ,
\end{equation}
the process is called invertible when the roots of $\theta(z) = 0$ lie outside
the unit circle. Invertibility means the model can be rewritten with the
unobservable shocks replaced by observable past values, as an AR process of
infinite order. This matters in practice because the $\epsilon_t$ are never
observed directly: they must be reconstructed from the data, and only an
invertible model makes that reconstruction unique and stable.

\paragraph{Identification.} The autocorrelation of an MA($q$) is zero for all
$k > q$, since observations more than $q$ periods apart share no shocks. For
MA(1) in particular,
\begin{equation}
\rho_1 = \frac{\theta_1}{1+\theta_1^2},
\qquad
\rho_k = 0 \quad \text{for } k \ge 2 .
\end{equation}
The signature is thus the mirror image of the AR case, and this is exactly
what appears in Figure~\ref{fig:ma1}: an ACF that cuts off sharply with a PACF
that decays gradually indicates an MA process. For $q=1$ the relation above inverts in
closed form, but for $q>1$ the coefficients cannot in general be recovered
from the autocorrelations this way, so they are estimated by
maximum likelihood, in practice through the Kalman filter applied to a state
space representation of the model.

\begin{figure}[ht]
\centering
\includegraphics[width=\linewidth]{figures/fig_ma1.pdf}
\caption{A simulated MA(1) series with $\theta_1 = 0.8$ ($n = 400$). The ACF
is significant at lag 1 only, matching $\rho_1 = 0.8/(1+0.64) \approx 0.49$,
while the PACF decays gradually -- the exact reverse of the AR signature in
Figure~\ref{fig:ar1}.}
\label{fig:ma1}
\end{figure}

% ---------------------------------------------------------------------
\subsection{ARIMA Model}

Stock prices are not stationary: they wander, and an AR or MA model fitted to
raw prices would violate its own assumptions. The ARIMA model resolves this by
differencing the series first. The full specification, ARIMA($p,d,q$), is

\begin{equation}
\phi_p(B)\,(1-B)^{d} X_t = \theta_q(B)\,\epsilon_t .
\end{equation}

Reading the three components:
\begin{itemize}
\item $\phi_p(B)$ is the autoregressive polynomial of Section~2.3, contributing
      dependence on $p$ past values.
\item $(1-B)^d$ is the differencing operator applied $d$ times. Since
      $(1-B)X_t = X_t - X_{t-1}$, applying it once converts a price into a
      price change; $d$ is almost always $0$ or $1$ in financial work, and
      $d=1$ on log prices produces log returns.
\item $\theta_q(B)$ is the moving-average polynomial of Section~2.4,
      contributing dependence on $q$ past shocks.
\end{itemize}

The effect of the middle term is shown in Figure~\ref{fig:diff}. The level
series wanders with no fixed mean and the ADF test cannot reject a unit root;
after one difference the series fluctuates around zero with a stable spread
and the unit root is decisively rejected. It is the differenced series that is
modelled.

\begin{figure}[ht]
\centering
\includegraphics[width=\linewidth]{figures/fig_differencing.pdf}
\caption{Differencing in practice on a simulated price series ($n = 600$).
Above: the plotted level, which wanders without a fixed mean; the ADF test
applied to it returns $p = 0.65$ and cannot reject a unit root. Below: the
first difference of the log price, which fluctuates about zero with stable
dispersion and for which the unit root is rejected at any conventional level
($p < 10^{-20}$).}
\label{fig:diff}
\end{figure}

\paragraph{The Box--Jenkins procedure.} Fitting an ARIMA model is an iterative
three-stage cycle \cite{Box_2015}. \emph{Identification} determines $d$ from
unit-root tests and then $(p,q)$ from the ACF and PACF of the differenced
series, using the two signatures established above. \emph{Estimation} obtains
the coefficients by maximum likelihood. \emph{Diagnostic checking} tests
whether the residuals are consistent with white noise; if they are not, the
specification is revised and the cycle repeats.

\paragraph{Choosing between candidate orders.} Because a larger model always
fits the training sample better, orders are compared by the Akaike Information
Criterion,
\begin{equation}
\mathrm{AIC} = -2\ln \hat{L} + 2k ,
\end{equation}
in which $\hat{L}$ is the maximised likelihood, so that $-2\ln\hat{L}$
measures lack of fit, and $k$ is the number of estimated parameters, so that
$2k$ is the penalty for complexity. Minimising the sum trades one against the
other. In small samples the corrected form
\begin{equation}
\mathrm{AICc} = \mathrm{AIC} + \frac{2k(k+1)}{n-k-1}
\end{equation}
is preferred, where $n$ is the number of observations; the correction vanishes
as $n$ grows.

\paragraph{Testing the residuals.} The Ljung--Box statistic tests whether the
first $m$ residual autocorrelations are jointly zero \cite{Ljung_1978}:
\begin{equation}
Q(m) = n(n+2)\sum_{k=1}^{m}\frac{\hat\rho_k^{\,2}}{n-k}
\;\sim\; \chi^2_{m-k_{\mathrm{par}}} ,
\end{equation}
where $\hat\rho_k$ is the residual autocorrelation at lag $k$, $n$ the sample
size, and $k_{\mathrm{par}}$ the number of fitted ARMA parameters, subtracted
from the degrees of freedom because those parameters were chosen using the same
data. A small $p$-value means structure remains in the residuals and the model
is inadequate.

\paragraph{Limitations.} ARIMA assumes both a linear conditional mean and a
constant variance. Neither describes financial returns well: returns show
little linear predictability but pronounced, persistent changes in volatility.
This motivates the seasonal extension in Section~2.6, the long-memory
extension in Section~2.9 and the volatility models in Section~2.10.

% ---------------------------------------------------------------------
\subsection{SARIMA Model}

Some series repeat on a fixed calendar cycle. SARIMA adds a second, seasonal
copy of the ARIMA structure operating at that cycle length:

\begin{equation}
\Phi_P(B^{s})\,\phi_p(B)\,(1-B^{s})^{D}(1-B)^{d}X_t
= \Theta_Q(B^{s})\,\theta_q(B)\,\epsilon_t .
\end{equation}

The equation looks formidable but is simply the ARIMA equation multiplied by
its own seasonal analogue:
\begin{itemize}
\item $s$ is the length of the cycle: $5$ for a weekly pattern in daily
      trading data, $12$ for a monthly pattern in monthly data.
\item $B^{s}$ is the backshift operator applied $s$ times, so that
      $B^{s}X_t = X_{t-s}$ reaches back exactly one full cycle.
\item $\Phi_P(B^{s})$ and $\Theta_Q(B^{s})$ are the seasonal autoregressive
      and moving-average polynomials. They relate this January to last
      January, rather than to last December.
\item $(1-B^{s})^{D}$ is seasonal differencing, subtracting the value one
      cycle earlier to remove a repeating pattern, just as $(1-B)^d$ removes a
      trend.
\end{itemize}

Seasonality is identified from the correlogram, where it produces spikes at
lags that are multiples of $s$, as in Figure~\ref{fig:seasonal}. For daily
equity prices such deterministic cycles are rarely strong, which is itself
consistent with weak-form efficiency: a reliable day-of-week effect would be
trivially exploitable and would therefore be arbitraged away. SARIMA is
retained here as the appropriate benchmark for the auxiliary series that do
show genuine periodicity, such as trading volume or macroeconomic covariates.

\begin{figure}[ht]
\centering
\includegraphics[width=\linewidth]{figures/fig_seasonal.pdf}
\caption{A simulated series with a twelve-period cycle. The correlogram shows
pronounced spikes at $k = 12$ and $k = 24$, the fingerprint of seasonality
that the $B^{s}$ terms in a SARIMA model are designed to absorb.}
\label{fig:seasonal}
\end{figure}

% ---------------------------------------------------------------------
\subsection{Exponential Smoothing}

Exponential smoothing reaches a forecast from a different direction. Rather
than modelling correlations, it maintains a running estimate of the current
level of the series and updates it as each new observation arrives. In plain
terms it is a weighted average in which recent observations count more and
older ones fade smoothly into irrelevance. Simple exponential smoothing, for a
series with neither trend nor seasonality, is

\begin{equation}
\ell_t = \alpha X_t + (1-\alpha)\,\ell_{t-1},
\qquad 0 < \alpha < 1,
\end{equation}
with the forecast for every future period given by
$\hat{X}_{t+h\mid t} = \ell_t$.

Term by term:
\begin{itemize}
\item $\ell_t$ is the estimated level of the series after seeing observation
      $t$. It is not observed in the data; it is a state the model carries
      forward.
\item $\alpha$ is the smoothing parameter, the fraction of the new
      observation absorbed into the level. Values near $1$ make the level
      track the data closely and react fast; values near $0$ make it smooth
      and slow.
\item $(1-\alpha)\ell_{t-1}$ is what is retained of the previous level.
\end{itemize}

\paragraph{Why the weights are exponential.} Substituting the recursion into
itself repeatedly gives
\begin{equation}
\ell_t = \alpha \sum_{k=0}^{t-1}(1-\alpha)^{k} X_{t-k}
        + (1-\alpha)^{t}\ell_0 .
\end{equation}
The weight on an observation $k$ periods old is $\alpha(1-\alpha)^k$, which
decays geometrically in $k$ -- the left panel of Figure~\ref{fig:ets}. The
weights are always positive and sum to one, so the level is a genuine weighted
average, and the influence of the arbitrary starting value $\ell_0$ decays
geometrically as well.

\begin{figure}[ht]
\centering
\includegraphics[width=\linewidth]{figures/fig_ets.pdf}
\caption{Exponential smoothing. Left: the weight $\alpha(1-\alpha)^k$ placed
on an observation $k$ periods old, for three values of $\alpha$. Right: the
resulting level $\ell_t$ on a simulated noisy series; $\alpha = 0.1$ filters
almost all of the noise but lags genuine movements, while $\alpha = 0.6$
follows the data closely at the cost of tracking the noise.}
\label{fig:ets}
\end{figure}

\paragraph{Adding a trend.} Holt's method carries a second state, the slope,
updated by the same idea:
\begin{equation}
\ell_t = \alpha X_t + (1-\alpha)(\ell_{t-1}+b_{t-1}),
\qquad
b_t = \beta^{*}(\ell_t - \ell_{t-1}) + (1-\beta^{*})\,b_{t-1},
\end{equation}
with the $h$-step forecast
\begin{equation}
\hat{X}_{t+h\mid t} = \ell_t + h\,b_t .
\end{equation}
Here $b_t$ estimates the current per-period increment and $\beta^{*}$ is its
own smoothing parameter: the new evidence about the slope is the change in
level $\ell_t - \ell_{t-1}$, and $\beta^{*}$ decides how much of it to accept.
The Holt--Winters extension adds a third state for a seasonal pattern.

\paragraph{Estimation and relationship to ARIMA.} The smoothing parameters and
initial states are estimated by maximum likelihood within the ETS (Error,
Trend, Seasonal) state space framework, with the particular variant chosen by
AIC exactly as in Section~2.5 \cite{49018680cf534023b6ae0f4e933c475d}. Some
ETS models coincide with ARIMA models -- simple exponential smoothing is
equivalent to ARIMA$(0,1,1)$, and Holt's linear method to ARIMA$(0,2,2)$ --
but the two families are parameterised differently, are selected by different
procedures and need not agree in finite samples, which is why both are kept as
separate benchmarks here.

% ---------------------------------------------------------------------
\subsection{Fractional Calculus Background}

The models so far allow memory that fades geometrically (AR), stops abruptly
(MA), or never fades at all (the random walk). Financial data, particularly
volatility, often behave in between: dependence that persists for months yet
still decays. Representing this requires differencing by a non-integer amount,
which in turn requires generalising differentiation itself.

Fractional calculus extends classical integer-order differentiation and
integration to arbitrary real orders. Benchohra et al.\ present several such
operators, including generalised fractional integrals, Riemann--Liouville
fractional derivatives, Caputo derivatives and Hilfer-type fractional
derivatives \cite{Benchohra_2023}. As an example, the $\psi$-Hilfer fractional
derivative of order $\alpha$ and type $\beta$, with $n-1 < \alpha < n$ and
$0 \le \beta \le 1$, is defined as

\begin{equation}
\left({}^{H}D_{a+}^{\alpha,\beta;\psi} h\right)(t) =
\left(I_{a+}^{\beta(n-\alpha);\psi}
\left(\frac{1}{\psi'(t)}\frac{d}{dt}\right)^{n}
I_{a+}^{(1-\beta)(n-\alpha);\psi} h\right)(t).
\end{equation}

The structure is what matters here rather than the analytical detail: a
fractional derivative is assembled by sandwiching ordinary differentiation,
the middle factor, between two fractional integrations, the operators
$I_{a+}$. The order $\alpha$ controls how much differentiation is applied and
the type $\beta$ selects the variant, so that the Hilfer derivative
interpolates between the Riemann--Liouville and Caputo derivatives as special
cases.

The relevance to time-series modelling is structural. ARIMA differences an
integer number of times, applying $(1-B)^d$ with $d \in \{0,1,2,\dots\}$;
fractional differencing applies the same operator with $d$ real. Expanding it
by the generalised binomial theorem gives
\begin{equation}
(1-B)^{d} X_t = \sum_{k=0}^{\infty}\binom{d}{k}(-1)^{k} X_{t-k},
\qquad
\binom{d}{k} = \frac{\Gamma(d+1)}{\Gamma(k+1)\,\Gamma(d-k+1)},
\end{equation}
where $\Gamma(\cdot)$ is the gamma function, the extension of the factorial to
non-integer arguments. In computation the equivalent recursion
\begin{equation}
\pi_0 = 1, \qquad
\pi_k = \pi_{k-1}\,\frac{k-1-d}{k},
\qquad
(1-B)^{d}X_t = \sum_{k=0}^{\infty}\pi_k X_{t-k},
\end{equation}
is preferred, since it avoids evaluating $\Gamma$ near its poles at the
non-positive integers. The recursion also makes the two regimes transparent.
When $d$ is a positive integer the factor $k-1-d$ is exactly zero at $k=d+1$,
so every subsequent coefficient vanishes and the sum terminates, recovering
ordinary differencing: for $d=1$ it collapses to $X_t - X_{t-1}$. When $d$ is
not an integer the factor is never zero, no coefficient ever vanishes, the sum
runs forever, and the present value depends on the entire history with weights
that decay slowly.
That infinite, slowly-decaying dependence is precisely the long memory
property exploited in the next subsection.

% ---------------------------------------------------------------------
\subsection{ARFIMA Model}

ARFIMA($p,d,q$) is the ARIMA model with the differencing parameter freed to
take real values:

\begin{equation}
\phi_p(B)\,(1-B)^{d} X_t = \theta_q(B)\,\epsilon_t,
\qquad d \in \mathbb{R}.
\end{equation}

The parameter $d$ now measures the strength of memory on a continuum:
\begin{itemize}
\item $d = 0$ gives a standard ARMA process with short memory.
\item $0 < d < 0.5$ gives long memory: the process is stationary, but
      dependence decays slowly.
\item $d \ge 0.5$ gives a nonstationary process, and $d = 1$ returns the
      integrated case of ARIMA.
\item $-0.5 < d < 0$ gives so-called intermediate memory, or antipersistence.
\end{itemize}
The process is stationary and invertible for $-0.5 < d < 0.5$. In that range
the autocorrelation decays as a power of the lag rather than geometrically,
\begin{equation}
\rho_k \sim C\,k^{2d-1} \quad \text{as } k \to \infty,
\end{equation}
where $C$ is a constant. The practical difference is large: a power-law decay
at lag 100 can still be visibly positive, whereas an AR(1) with $\phi_1 = 0.7$
has $\rho_{100} = 0.7^{100} \approx 3\times10^{-16}$. The parameter is often
reported through the Hurst exponent
\begin{equation}
H = d + 0.5 ,
\end{equation}
so that $H = 0.5$ corresponds to no long memory and $H > 0.5$ to persistence.

\paragraph{Estimation.} Exact maximum likelihood is expensive because the
infinite expansion must be truncated. Two practical alternatives are standard:
the semiparametric Geweke--Porter-Hudak estimator, which regresses the log
periodogram of the series on a log-frequency term at low frequencies and reads
$d$ off the slope; and the Whittle likelihood estimator, which is fully
parametric and asymptotically more efficient \cite{Granger_1980}.

\paragraph{Where long memory actually appears.} Figure~\ref{fig:longmem}
illustrates the model and simultaneously delivers one of the most important
facts in this section. The series is constructed so that its log-volatility is
a fractionally integrated process with $d = 0.35$; the returns are that
volatility multiplied by independent noise. The left panel shows the
correlogram of the returns themselves: essentially every autocorrelation lies
inside the white-noise band, so the returns are close to unforecastable,
exactly as weak-form efficiency predicts. The right panel shows the
correlogram of their \emph{absolute values}, which stays positive across a
hundred lags and follows the power law $k^{2d-1}$ implied by the construction.
The direction of tomorrow's move is nearly unpredictable; its likely
\emph{magnitude} is highly predictable far into the future.

It is worth being precise about what distinguishes this from the GARCH models
of the next subsection. A GARCH process with $\alpha_1+\beta_1$ close to one
also produces volatility that stays elevated for a long time, but its
autocorrelations still decay \emph{geometrically} and so eventually collapse;
long memory in the ARFIMA sense means \emph{hyperbolic} decay, which does not.
Distinguishing high persistence from true long memory is an empirical question
about the value of $d$, not something that can be settled by eye. ARFIMA is
therefore applied in this study to volatility proxies such as absolute or
squared returns rather than to prices, and the asymmetry between unpredictable
signs and predictable magnitudes is the direct motivation for the volatility
models that follow.

\begin{figure}[ht]
\centering
\includegraphics[width=\linewidth]{figures/fig_long_memory.pdf}
\caption{Two correlograms of the same simulated return series ($n = 4000$),
constructed so that its log-volatility is a fractionally integrated process
with $d = 0.35$. Left: the returns $r_t$, whose autocorrelations lie almost
entirely inside the 95\% white-noise band, indicating near-unpredictability of
direction. Right: the absolute returns $|r_t|$, which remain positive across a
hundred lags and track the theoretical power law $C k^{2d-1}$ (dashed) rather
than decaying geometrically.}
\label{fig:longmem}
\end{figure}

% ---------------------------------------------------------------------
\subsection{ARCH and GARCH Models}

Every model so far has assumed that the size of a typical shock is the same in
every period. Financial returns violate this conspicuously: turbulent days
cluster with turbulent days and quiet days with quiet days, as the upper panel
of Figure~\ref{fig:garch} shows. This phenomenon is called volatility
clustering, and modelling it requires letting the variance itself depend on the
past.

\paragraph{The ARCH model.} Engle's formulation separates the return into a
mean and a scaled shock \cite{Engle_1982}:
\begin{equation}
r_t = \mu + \epsilon_t ,
\qquad
\epsilon_t = z_t \sigma_t ,
\qquad
z_t \sim \mathrm{iid}(0,1),
\end{equation}
where $\sigma_t$ is the conditional standard deviation, the typical size of a
move given everything known up to $t-1$, and $z_t$ is a standardised random
draw that supplies only the sign and relative size. The variance then follows
\begin{equation}
\sigma_t^{2} = \alpha_0 + \sum_{i=1}^{q}\alpha_i \epsilon_{t-i}^{2},
\qquad \alpha_0 > 0,\; \alpha_i \ge 0 .
\end{equation}
Term by term:
\begin{itemize}
\item $\sigma_t^2$ is today's variance, forecast rather than assumed.
\item $\alpha_0$ is a floor: the baseline variance that remains even after a
      long calm stretch.
\item $\epsilon_{t-i}^{2}$ is the squared shock $i$ days ago. Squaring
      discards the sign, encoding the assumption that a large move in either
      direction raises tomorrow's expected volatility.
\item $\alpha_i$ measures how strongly that past shock feeds through. The
      non-negativity constraints exist because a variance cannot be negative.
\end{itemize}
Covariance stationarity requires $\sum_{i=1}^{q}\alpha_i < 1$.

\paragraph{The GARCH generalisation.} Capturing the long persistence of
volatility with ARCH alone requires many lags. Bollerslev's insight was to let
the variance depend on its own past as well \cite{Bollerslev_1986}:
\begin{equation}
\sigma_t^{2} = \alpha_0
+ \sum_{i=1}^{q}\alpha_i \epsilon_{t-i}^{2}
+ \sum_{j=1}^{p}\beta_j \sigma_{t-j}^{2}.
\end{equation}
The new term $\beta_j \sigma_{t-j}^2$ carries forward yesterday's variance
directly, so a single lag of each type reproduces very persistent volatility
with three parameters. The GARCH(1,1) model,
\begin{equation}
\sigma_t^{2} = \alpha_0 + \alpha_1 \epsilon_{t-1}^{2} + \beta_1 \sigma_{t-1}^{2},
\end{equation}
is the workhorse of applied volatility modelling. Here $\alpha_1$ governs how
sharply volatility reacts to a new shock and $\beta_1$ how long that elevated
volatility persists; empirically $\alpha_1$ is small and $\beta_1$ close to
one. The process is covariance-stationary, with finite unconditional variance
$\alpha_0/(1-\alpha_1-\beta_1)$, if and only if
\begin{equation}
\sum_{i=1}^{q}\alpha_i + \sum_{j=1}^{p}\beta_j < 1 .
\end{equation}
Figure~\ref{fig:garch} shows a GARCH(1,1) fitted to a simulated series whose
true conditional volatility is known; the estimated path tracks the truth
closely, recovering a quantity that is never directly observed.

\begin{figure}[ht]
\centering
\includegraphics[width=\linewidth]{figures/fig_garch.pdf}
\caption{Volatility clustering and its estimation. Above: returns simulated
from a GARCH(1,1) process, showing alternating calm and turbulent stretches.
Below: the true conditional volatility $\sigma_t$ used in the simulation
against the path $\hat\sigma_t$ estimated by a GARCH(1,1) fitted to the
returns alone.}
\label{fig:garch}
\end{figure}

\paragraph{Testing and estimation.} Whether a volatility model is needed at
all is checked with Engle's ARCH-LM test, which regresses squared residuals
from the mean equation on their own lags and tests whether those coefficients
are jointly zero. Parameters are estimated by maximum likelihood, commonly
with Student-$t$ rather than normal innovations, since returns have heavier
tails than the normal distribution allows even after conditioning on
$\sigma_t$.

\paragraph{Asymmetry.} Because $\epsilon_{t-1}$ enters squared, GARCH treats a
large fall and an equally large rise identically. In equity markets they are
not: falls raise subsequent volatility more, an effect known as the leverage
effect. The GJR-GARCH model adds a term active only for negative shocks,
\begin{equation}
\sigma_t^{2} = \alpha_0 + \alpha_1\epsilon_{t-1}^{2}
+ \gamma\,\epsilon_{t-1}^{2}\,\mathbf{1}_{\{\epsilon_{t-1}<0\}}
+ \beta_1\sigma_{t-1}^{2},
\end{equation}
in which the indicator $\mathbf{1}_{\{\epsilon_{t-1}<0\}}$ equals one after a
negative shock and zero otherwise, so a positive $\gamma$ makes bad news raise
volatility by more than good news of the same magnitude. The EGARCH model of
Nelson \cite{Nelson_1991} achieves the same asymmetry by modelling
$\ln \sigma_t^2$, which additionally removes the need for non-negativity
constraints.

% ---------------------------------------------------------------------
\subsection{Markov-Switching Models}

Every model above assumes one fixed set of parameters governs the entire
sample. Markets are widely observed to alternate instead between distinct
states -- a calm regime of small moves and mild upward drift, and a turbulent
regime of large moves and negative drift. A single fixed-coefficient model
must average across both, describing neither.

\paragraph{The model.} Following Hamilton \cite{Hamilton_1989}, let $S_t$ be
an unobserved state variable taking one of $K$ values, which evolves as a
first-order Markov chain:
\begin{equation}
p_{ij} = P(S_t = j \mid S_{t-1} = i),
\qquad \sum_{j=1}^{K} p_{ij} = 1 .
\end{equation}
The transition probability $p_{ij}$ is the chance of moving from state $i$ to
state $j$ in one period, and "first-order" means only the current state
matters for where the process goes next. The diagonal entries $p_{ii}$ measure
persistence: a value of $0.98$ implies an average stay of $1/(1-0.98) = 50$
periods in that regime. Conditional on the state, returns follow a
state-specific distribution; the standard specification in finance switches
both mean and variance,
\begin{equation}
r_t \mid S_t = k \;\sim\; N\!\left(\mu_k,\, \sigma_k^{2}\right),
\qquad k = 1,\dots,K .
\end{equation}
With $K=2$ the parameters have a direct reading: $\mu_1 > 0$ with small
$\sigma_1$ describes the calm regime, $\mu_2 < 0$ with large $\sigma_2$ the
turbulent one.

\paragraph{Estimation with an unobserved state.} The regime is never observed,
so the likelihood must average over all possible regime paths. The Hamilton
filter does this recursively, combining the transition probabilities with the
likelihood of each new observation to produce the filtered probability
\begin{equation}
P(S_t = k \mid X_1,\dots,X_t)
= \frac{f(X_t \mid S_t = k)\,P(S_t = k \mid X_1,\dots,X_{t-1})}
       {\sum_{j=1}^{K} f(X_t \mid S_t = j)\,P(S_t = j \mid X_1,\dots,X_{t-1})},
\end{equation}
where $f(X_t \mid S_t = k)$ is the density of the observation under regime
$k$. The numerator weights each regime by how well it explains today's
observation, and the denominator normalises so the probabilities sum to one.
Maximising the resulting likelihood, in practice via the EM algorithm, yields
the parameters. Smoothed probabilities $P(S_t = k \mid X_1,\dots,X_T)$, which
use the whole sample, are reported for interpretation, but only the filtered
probabilities are admissible as trading signals, since the smoothed version
uses future data and would violate the constraint in Section~2.1.

Figure~\ref{fig:markov} shows a two-regime model fitted to a simulated series.
The model, given only the returns, recovers the regime dates almost exactly.

\begin{figure}[ht]
\centering
\includegraphics[width=\linewidth]{figures/fig_markov.pdf}
\caption{A two-regime Markov-switching model on a simulated series ($n=900$).
Above: the returns, with the periods truly generated by the high-variance
regime shaded. Below: the \emph{filtered} probability of being in that
regime, that is, the estimate available in real time using only data up to
each date. The fitted regime standard deviations,
$0.71$ and $2.39$, closely recover the values $0.7$ and $2.3$ used to
generate the data.}
\label{fig:markov}
\end{figure}

\paragraph{Positioning.} Markov-switching models occupy a useful middle
ground. They remain fully parametric and likelihood-based, like everything
else in this section, so they retain interpretability and formal inference;
yet by allowing the parameters themselves to change they capture a genuine
form of nonlinearity that no ARIMA or GARCH specification can. This makes them
the natural bridge to the machine learning methods of the next section.

% ---------------------------------------------------------------------
\subsection{Transition to Machine Learning}

Taken together, the models of this section share a common architecture. Each
specifies in advance a functional form -- linear in past values, linear in past
shocks, linear in past squared shocks, or linear within each of a small number
of regimes -- and then estimates a modest number of parameters within that
form. This is the source of both their strength and their limitation. Because
the form is fixed, the parameters are interpretable, the estimators have known
asymptotic properties, and there are formal tests for when the assumptions
fail. But if the true relationship between past and future is not of the
assumed form, no amount of data will reveal it.

The models of the following section reverse this trade. They impose almost no
structure on the functional form, learning it from the data instead, and in
exchange give up closed-form inference and much of the interpretability. The
comparison this paper undertakes is between these two positions, conducted
under the single evaluation protocol fixed in Section~2.1 and measured against
the random-walk benchmark of Section~2.2.

\bibliographystyle{plain}
\bibliography{refs}

\end{document}